\chapter{第一题}

\section{题目}

死锁是并行程序中的典型错误之一，以下代码段实现了在两个进程上交换各自长度为 N 的浮点数组的功能。
\begin{lstlisting}[style=cleancpp]
float send_buf[N], recv_buf[N];
if(my_rank== 0){ //进程 0 
    MPI_Send(send_buf, N, MPI_FLOAT, 1, 0, MPI_COMM_WORLD); 
    MPI_Recv(recv_buf, N, MPI_FLOAT, 1, 1, MPI_COMM_WORLD, MPI_STATUS_IGNORE); 
} else if(my_rank== 1){ //进程 1 
    MPI_Send(send_buf, N, MPI_FLOAT, 0, 1, MPI_COMM_WORLD); 
    MPI_Recv(recv_buf, N, MPI_FLOAT, 0, 0, MPI_COMM_WORLD, MPI_STATUS_IGNORE); 
}
\end{lstlisting}

\begin{enumerate}
    \item 分析该程序为什么在 N 较小时能够成功运行，而当 N 的取值非常大时，有可能永久挂起。
    \item 请给出一种可行的解决方案并简要说明。
\end{enumerate}

\section{答案}
\subsection{（1）原因分析} 

当两个进程都先调用阻塞式发送 \texttt{MPI\_Send} 而未先发布对应的接收 \texttt{MPI\_Recv} 时，若消息大小 $N$ 超过 MPI 实现内部的 "eager" 缓冲阈值，发送操作会阻塞等待对端接收；双方都在 \texttt{MPI\_Send} 上等待，从而发生死锁。对于较小的 $N$，MPI 常使用 eager 协议把数据复制到内部缓冲区使发送立即返回，因此看起来没有问题；但当 $N$ 很大时，发送不会被缓冲而会阻塞，进而可能导致永久挂起。

\subsection{（2）可行的解决方案}
死锁的根因是阻塞发送与缓冲策略，解决方法是使用可避免互相等待的通信模式（\texttt{MPI\_Sendrecv}、非阻塞或先贴接收）。


\subsubsection{使用原子式的发送/接收函数 \texttt{MPI\_Sendrecv}：}

该调用在内部避免了双方同时阻塞，能可靠交换任意大小的数据（由实现保证）。

\begin{lstlisting}[style=cleancpp]
int peer = 1 - my_rank;
MPI_Sendrecv(send_buf, N, MPI_FLOAT, peer, 0, 
             recv_buf, N, MPI_FLOAT, peer, 0,
             MPI_COMM_WORLD, MPI_STATUS_IGNORE);
\end{lstlisting}

\subsubsection{使用非阻塞通信先贴接收再发送：}

先发布接收可以保证对端的发送不会无限等待，非阻塞发送/接收配合等待也能避免死锁。

\begin{lstlisting}[style=cleancpp]
MPI_Request reqs[2];
MPI_Irecv(recv_buf, N, MPI_FLOAT, peer, 0, MPI_COMM_WORLD, &reqs[0]);
MPI_Isend(send_buf, N, MPI_FLOAT, peer, 0, MPI_COMM_WORLD, &reqs[1]);
MPI_Waitall(2, reqs, MPI_STATUSES_IGNORE);
\end{lstlisting}
    
\subsubsection{改变顺序：}

一端先执行 \texttt{MPI\_Recv} 再 \texttt{MPI\_Send}，另一端相反。或者，使用 \texttt{MPI\_Isend}/\texttt{MPI\_Irecv} 的混合策略以保证不存在循环等待。

\chapter{第二题}

\section{题目}

在分布式算法中，有时我们需要实现类似于 \texttt{MPI\_Allgather} 的功能，即每个进程都有一个局部值，最后每个进程都要获得所有进程的值，实现全局数据同步。使用 C 语言编写一段符合以下要求的程序：
\begin{enumerate}
    \item 在一个包含 \texttt{size} 个进程的环形传递结构（即进程 q 总是将消息发送给 \texttt{(q+1) mod P}，并从 \texttt{(q-1+P) mod P} 接收消息，其中 P 为进程总数）中，实现 2 中的要求。
    \item 每个进程持有一个整数 \texttt{my\_val}。经过若干步迭代，每个进程计算出环内所有值的累加和。
    \item 使用 \texttt{MPI\_Sendrecv} 以避免死锁。
\end{enumerate}

\section{答案}
下面给出一种在环形结构上使用 \texttt{MPI\_Sendrecv} 的实现思路和示例代码：每个进程将自己的值沿环传递，经过 $P-1$ 步后，每个进程都能收到环上所有进程的值并累加。

\begin{lstlisting}[style=cleancpp]
#include <mpi.h>
#include <stdio.h>

int main(int argc, char **argv) {
    MPI_Init(&argc, &argv);
    int rank, size;
    MPI_Comm_rank(MPI_COMM_WORLD, &rank);
    MPI_Comm_size(MPI_COMM_WORLD, &size);

    int my_val; 
    my_val = rank + 1;

    int sum = my_val;
    int send_val = my_val;
    int recv_val;
    int right = (rank + 1) % size;
    int left  = (rank - 1 + size) % size;

    for (int step = 1; step < size; ++step) {
        MPI_Sendrecv(&send_val, 1, MPI_INT, right, 0,
                     &recv_val, 1, MPI_INT, left,  0,
                     MPI_COMM_WORLD, MPI_STATUS_IGNORE);
        sum += recv_val;
        send_val = recv_val;
    }

    printf("rank %d: sum = %d\n", rank, sum);

    MPI_Finalize();
    return 0;
}
\end{lstlisting}

说明：每个进程在第 $i$ 步发送上一步收到的值（第 0 步发送自己的值），经过 $P-1$ 步后，环上所有值都已被遍历并累加。使用 \texttt{MPI\_Sendrecv} 可以在一次调用内完成发送和接收，避免了阻塞发送引起的死锁。

\chapter{第三题}

\section{题目}

有两个长度为 $(n=12)$ 的整数向量 x 和 y，其初值如下：
$$ x=[1, 3, 5, 7, 9, 11, 13, 15, 17, 19, 21, 23] $$
$$ y=[2, 4, 6, 8, 10, 12, 14, 16, 18, 20, 22, 24] $$
现使用 \texttt{comm\_sz= 4} 个 MPI 进程完成向量加法，要求如下：
\begin{enumerate}
    \item 由 0 号进程初始化向量 x 和 y；
    \item 使用 MPI 通信操作，将两个向量平均分发给各个进程；
    \item 每个进程计算自己负责部分的局部和；
    \item 再将结果收集到 0 号进程；
    \item 最后由 0 号进程输出结果向量 z。
\end{enumerate}

请回答下列问题：
\begin{enumerate}
    \item 每个进程分别处理向量中的哪几个元素？
    \item 最终结果向量 z 是什么？
    \item 写出完成该任务时需要用到的主要 MPI 函数名称，并简要说明它们的作用。
\end{enumerate}

\section{答案}
\begin{enumerate}
    \item 每个进程分别处理的元素（下标按从 0 开始计数，向量长度 $n=12$，进程数 $P=4$，每个进程处理 $n/P=3$ 个元素）：
    \begin{itemize}
        \item 进程 0: 元素下标 $0,1,2$ （对应 $x=[1,3,5]$, $y=[2,4,6]$）
        \item 进程 1: 元素下标 $3,4,5$ （对应 $x=[7,9,11]$, $y=[8,10,12]$）
        \item 进程 2: 元素下标 $6,7,8$ （对应 $x=[13,15,17]$, $y=[14,16,18]$）
        \item 进程 3: 元素下标 $9,10,11$（对应 $x=[19,21,23]$, $y=[20,22,24]$）
    \end{itemize}

    \item 最终结果向量 $z$（逐元素相加）为：
    $$z=[3,7,11,15,19,23,27,31,35,39,43,47].$$

    \item 完成该任务常用的主要 MPI 函数及作用：
    \begin{itemize}
        \item \texttt{MPI\_Init}: 初始化 MPI 环境。
        \item \texttt{MPI\_Comm\_size}: 获取进程总数（$P$）。
        \item \texttt{MPI\_Comm\_rank}: 获取当前进程的秩（rank）。
        \item \texttt{MPI\_Scatter}: 将根进程上的数组按等份分发到各个进程（每个进程接收相同数量的数据）。
        \item \texttt{MPI\_Gather}: 将各进程计算出的局部结果收集到根进程。
        \item \texttt{MPI\_Finalize}: 终止 MPI 环境。
    \end{itemize}
若数据块大小不一致，应使用 \texttt{MPI\_Scatterv}/\texttt{MPI\_Gatherv} 或手工发送/接收。
\end{enumerate}

\chapter{第四题}

\section{题目}

设有一个长度为 $(n=10)$ 的整数数组 a：
$$ a=[5, 8, 2, 6, 1, 9, 4, 7, 3, 10] $$
现在希望使用 \texttt{comm\_sz= 4} 个 MPI 进程对该数组中的元素分别乘以 2，最后再把结果收集回 0 号进程。

已知 \texttt{MPI\_Scatter} 和 \texttt{MPI\_Gather} 的一个限制是：每个进程接收或发送的数据个数必须相同。

请回答下列问题：
\begin{enumerate}
    \item 为什么在本题中不能直接使用一次 \texttt{MPI\_Scatter} 和一次 \texttt{MPI\_Gather} 完成任务？
    \item 如果采用“前 2 个进程各处理 3 个元素，后 2 个进程各处理 2 个元素”的分配方案，请写出各进程负责的数据范围。
    \item 给出一种可行的修改思路，使程序在 $n$ 不能被 \texttt{comm\_sz} 整除时仍能正确完成计算。
\end{enumerate}

\section{答案}
\begin{enumerate}
    \item 不能直接使用一次 \texttt{MPI\_Scatter} 和一次 \texttt{MPI\_Gather} 的原因：这两个函数要求每个进程发送/接收的数据个数相同，但本题 $n=10$、进程数 $P=4$ 时 $10$ 不能被 $4$ 整除，无法在一次等份划分中保证每个进程接收相同数量的数据。因此必须使用变长的分发/收集或额外处理。

    \item 若采用“前 2 个进程各处理 3 个元素，后 2 个进程各处理 2 个元素”的分配方案，则各进程负责的数据范围（按 0 起始下标）：
    \begin{itemize}
        \item 进程 0: 下标 $0,1,2$（对应元素 $[5,8,2]$）
        \item 进程 1: 下标 $3,4,5$（对应元素 $[6,1,9]$）
        \item 进程 2: 下标 $6,7$  （对应元素 $[4,7]$）
        \item 进程 3: 下标 $8,9$  （对应元素 $[3,10]$）
    \end{itemize}

    \item 根进程（例如 rank 0）准备原始数组 \texttt{sendbuf}，并计算好 \texttt{counts} 与 \texttt{displs}；每个进程根据 \texttt{counts[rank]} 分配接收缓冲区 \texttt{recvbuf}，调用 \texttt{MPI\_Scatterv} 分发局部数据，做本地计算后再用 \texttt{MPI\_Gatherv} 收集回根进程。对任意长度和进程数，可用如下方法计算 \texttt{counts} 和 \texttt{displs}：
    \begin{itemize}
        \item 令 $q=\lfloor n/P\rfloor,\; r=n\bmod P$，则对每个进程 $i$，\texttt{counts[i]=q+(i<r?1:0)}。
        \item \texttt{displs[0]=0;} 对于 $i>0$，\texttt{displs[i]=displs[i-1]+counts[i-1];}
    \end{itemize}

    示例代码片段：
    \begin{lstlisting}[style=cleancpp]
int counts[4] = {3,3,2,2};
int displs[4] = {0,3,6,8};
int rank;
MPI_Comm_rank(MPI_COMM_WORLD, &rank);
int recvcount = counts[rank];
int *recvbuf = malloc(recvcount * sizeof(int));
// 根进程上有 int sendbuf[n]; 其它进程该参数可置为 NULL 
MPI_Scatterv(sendbuf, counts, displs, MPI_INT,
             recvbuf, recvcount, MPI_INT, 0, MPI_COMM_WORLD);
// 本地计算：每个元素乘以 2 
for (int i = 0; i < recvcount; ++i) recvbuf[i] *= 2;
// 收集回根进程 
MPI_Gatherv(recvbuf, recvcount, MPI_INT,
            result, counts, displs, MPI_INT, 0, MPI_COMM_WORLD);
free(recvbuf);
    \end{lstlisting}
\texttt{MPI\_Scatterv} 和 \texttt{MPI\_Gatherv} 在根进程上读取 \texttt{counts} 与 \texttt{displs}，非根进程这两个数组参数会被忽略；但各进程必须知道自己的接收长度（\texttt{recvcount}）以便正确分配接收缓冲区。该方法通用且避免了对数据填充或多次点对点发送的手工实现。
\end{enumerate}

\chapter{第五题}

\section{题目}

考虑使用 MPI 实现数值积分的梯形法则，假设函数 $f(x)$ 在区间 $[a，b]$ 上是连续的，我们将区间划分成 $n$ 个等宽的小梯形，每个梯形的宽度为 $\hbar = \frac{b-a}{n}$，积分的近似公式为：
$$ \int_a^b f(x)dx \approx \hbar \left[ \frac{f(a)}{2} + \sum_{i=1}^{n-1} f(a+ i \ast \hbar) + \frac{f(b)}{2} \right] $$
现在你编写了一个 MPI 程序，使用 $p$ 个进程来并行计算该积分，每个进程负责计算一部分梯形的面积，然后通过 \texttt{MPI\_Reduce} 将局部结果汇总到进程 0，最后由进程 0 输出总积分值。

\begin{enumerate}
    \item 假设 $n$ 能被 $p$ 整除，请写出每个进程负责的梯形数量、起始点和结束点的计算公式。如果 $n$ 不能被 $p$ 整除，每个进程负责的梯形数量又是多少。
    \item 使用 \texttt{MPI\_Reduce} 汇总结果，该集合通信的时间复杂度是多少，假设使用的是树形通信结构。
    \item 假设串行程序的运行时间为 $T_{serial}= k \ast n$，并行程序的运行时间为：$T_{parallel}(n, p)= \frac{k \ast n}{p} + c \ast \log_2(p)$, 其中 $c$ 为常数，请推导出加速比 $S(n, p)$ 和效率 $E(n, p)$ 的表达式。当 $n$ 固定，$p$ 增大时，效率会如何变化，为什么。
\end{enumerate}

\section{答案}
\begin{enumerate}
    \item \textbf{当 $n$ 能被 $p$ 整除时：}每个进程负责的梯形数量为
    $$n_{local} = \frac{n}{p}.$$ 对于秩为 $\text{rank}$ 的进程，其负责的区间为
    $$x_{start} = a + \text{rank} \cdot n_{local} \cdot \hbar,\qquad x_{end} = a + (\text{rank}+1) \cdot n_{local} \cdot \hbar,$$
    其中 $\hbar=(b-a)/n$。
    
    \textbf{如果 $n$ 不能被 $p$ 整除：}
    可令\begin{align*}
        q &= \left\lfloor\frac{n}{p}\right\rfloor, & r &= n \bmod p, \\
        n_{local}(\text{rank}) &= q + (\text{rank}<r?1:0),
    \end{align*}
    即前 $r$ 个进程多分配 1 个梯形。

    \item 使用树形归约（树形通信结构）的 \texttt{MPI\_Reduce} 的时间复杂度为 $O(\log p)$。若要更精确地包含延迟项和带宽项，通常写作 $O(\alpha\log p + \beta m\log p)$。若只关心步数为 $\Theta(\log p)$。

    \item 给定
    $$T_{serial}=k\,n,\qquad T_{parallel}(n,p)=\frac{k\,n}{p}+c\,\log_2 p,$$
    则加速比为
    $$S(n,p)=\frac{T_{serial}}{T_{parallel}}=\frac{k\,n}{\dfrac{k\,n}{p}+c\,\log_2 p}=\frac{k\,n\,p}{k\,n+c\,p\,\log_2 p}.$$  效率为
    $$E(n,p)=\frac{S(n,p)}{p}=\frac{k\,n}{k\,n+c\,p\,\log_2 p}=\frac{1}{1+\dfrac{c\,p\,\log_2 p}{k\,n}}.$$
    当固定 $n$，$p$ 增大时，分母中通信项 $\dfrac{c\,p\,\log_2 p}{k\,n}$ 会增长，导致效率下降；直观上，通信开销随进程数增加（且近似为 $p\log p$ 级别）增长，最终通信会掩盖并行计算的收益，效率趋向 0。
\end{enumerate}